%============================================================
% CH1 · P4 · PROBLEM 1 — PLUS ONE (LeetCode 66)
%============================================================
\subsection{Problem 1.4.1: Plus One}
\problemlabel{1.4.1}{ssec:ch1-plus-one}

\problemcard{1.4.1}{Plus One}{LeetCode 66}{https://leetcode.com/problems/plus-one/}{\exeasy\ \lc{Easy}}{Simulate the carry from the last digit; stop at the first digit that is not 9; grow only for $99\ldots9$}

\subsubsection*{1. Problem statement}

A non-negative integer is given as an array \texttt{digits} of its decimal digits, \textbf{most significant digit first}, without leading zeros (except for the number 0 itself). Return the digits of the number plus one, in the same format.

\textbf{Constraints:} $1 \le n \le 100$, \; $0 \le \texttt{digits[i]} \le 9$, \; no leading zeros.

\begin{table}[H]
\centering\small
\begin{tabular}{@{}llll@{}}
\toprule
\textbf{digits} & \textbf{Output} & \textbf{Number} & \textbf{What happens} \\
\midrule
$[1,2,3]$ & $[1,2,4]$ & $123 + 1 = 124$ & last digit grows, no carry \\
$[4,3,2,1]$ & $[4,3,2,2]$ & $4321 + 1 = 4322$ & last digit grows, no carry \\
$[9]$ & $[1,0]$ & $9 + 1 = 10$ & carry out of the top digit: one digit longer \\
\bottomrule
\end{tabular}
\end{table}

\begin{diagrambox}
\begin{tikzpicture}[x=1cm, y=1cm]
    \begin{scope}
        \dparray{a}{0.8}{in}{1/acell, 2/acell, 3/aread}
        \dparray{b}{-0.4}{out}{1/acell, 2/acell, 4/awrite}
        \node[font=\scriptsize\bfseries] at (1.35,2.05) {Sample 1};
    \end{scope}
    \begin{scope}[xshift=4.6cm]
        \dparray{c}{0.8}{}{4/acell, 3/acell, 2/acell, 1/aread}
        \dparray{d}{-0.4}{}{4/acell, 3/acell, 2/acell, 2/awrite}
        \node[font=\scriptsize\bfseries] at (1.8,2.05) {Sample 2};
    \end{scope}
    \begin{scope}[xshift=9.6cm]
        \dparray{e}{0.8}{}{{}/emptycell, 9/aread}
        \dparray{f}{-0.4}{}{1/awrite, 0/awrite}
        \node[font=\scriptsize\bfseries] at (0.9,2.05) {Sample 3};
        \draw[dep] ([yshift=1mm]e1.north) to[bend right=50] node[lbl, above]{carry} ([yshift=1mm]e0.north);
    \end{scope}
\end{tikzpicture}
\figcap{The official examples. Orange: the digit that receives the $+1$; blue: digits written by the algorithm. In Sample 3 the array must grow by one cell at the front.}
\end{diagrambox}

\subsubsection*{2. The problem model}
\begin{modelbox}[Problem Model --- Plus One]
\renewcommand{\arraystretch}{1.35}
\begin{tabularx}{\linewidth}{@{}>{\bfseries}l X@{}}
Input & a digit array $d[0..n-1]$, $d[0]$ most significant; value $N = \sum_i d[i]\cdot 10^{\,n-1-i}$ \\
Computed & the digit array of $N + 1$ (length $n$ or $n+1$) \\
Objects & the digits, processed from the least significant end; the state is one carry bit \\
Size & $n \le 100$: $N$ can have 100 digits --- far beyond \texttt{long long} ($\approx 9.2\cdot10^{18}$, 19 digits); any $\bigO{n}$ digit loop is instant \\
\end{tabularx}
\end{modelbox}

\begin{diagrambox}
\begin{tikzpicture}[x=1cm, y=1cm]
    \dparray{m}{0}{digits}{$d_0$/acell, $\cdots$/acell, $d_{j-1}$/acell, $d_j$/aread, 9/aout, 9/aout, 9/aout}
    \arange{m4}{m6}{trailing 9s: all become 0}{memoorange}
    \aptrup{m3}{first non-9 from the right: $+1$}{accentorange}
    \node[note, anchor=west, text width=3.6cm, align=left] at (6.8,0.3) {digits left of $d_j$ are copied unchanged};
\end{tikzpicture}
\figcap{The whole structure of the answer: find the last digit that is not 9, add one to it, and zero the run of 9s behind it. If every digit is 9, the answer is $1$ followed by $n$ zeros.}
\end{diagrambox}

\subsubsection*{3. Key observations}

\begin{itemize}[leftmargin=*]
    \item \textbf{The number does not fit in any integer type.} With 100 digits, ``convert, add one, convert back'' overflows. Even \texttt{unsigned long long} wraps around after about $1.8 \cdot 10^{19}$: on twenty 9s the conversion approach returns \texttt{7766279631452241920} instead of $1$ followed by twenty zeros. We must simulate the addition on the digits.
    \item \textbf{The carry is 0 or 1, and it continues only through 9s.} $d + 1 \le 9$ for $d \le 8$: no carry, and nothing to the left changes. $9 + 1 = 10$: write 0, carry 1. So the loop can stop at the first digit (from the right) that is not 9.
    \item \textbf{The array grows only for $99\ldots9$.} If the carry leaves the most significant digit, every digit was 9 and is now 0; the result is $1\,0\,0 \ldots 0$ with $n+1$ digits. This is the single border case.
\end{itemize}

\subsubsection*{4. Approaches}

\begin{bruteforcebox}
\textbf{Idea.} Convert the digits into an integer, add one, convert back.
\begin{lstlisting}[style=cpp]
#include <bits/stdc++.h>
using namespace std;

class Solution {
public:
    vector<int> plusOne(vector<int>& digits) {
        unsigned long long x = 0;
        for (int d : digits) x = x * 10 + d;      // overflows beyond about 1.8e19
        x++;
        vector<int> res;
        do { res.push_back(x % 10); x /= 10; } while (x > 0);
        reverse(res.begin(), res.end());
        return res;
    }
};

int main() {                                      // input: n, then the n digits
    int n;
    if (!(cin >> n)) return 0;
    vector<int> d(n);
    for (auto& x : d) cin >> x;
    vector<int> r = Solution().plusOne(d);
    for (size_t i = 0; i < r.size(); i++) cout << r[i] << (i + 1 < r.size() ? ' ' : '\n');
}
\end{lstlisting}
\begin{complexitybox}
\textbf{Time:} $\bigO{n}$, \textbf{Space:} $\bigO{n}$ --- fast but \wa{}: correct for at most 19 digits, wrong on twenty 9s (it prints \texttt{7766279631452241920}). This is the typical ``brute force'' of digit problems: not slow, but limited by the integer type.
\end{complexitybox}
\end{bruteforcebox}

\begin{algorithmbox}
\textbf{Idea.} Walk from the last digit to the first. A digit below 9 is incremented and the answer is complete. A 9 becomes 0 and the carry moves one place to the left. If the loop runs off the front, insert a 1.
\begin{lstlisting}[style=cpp]
#include <bits/stdc++.h>
using namespace std;

class Solution {
public:
    vector<int> plusOne(vector<int>& digits) {
        for (int i = (int)digits.size() - 1; i >= 0; i--) {
            if (digits[i] < 9) {          // no carry out of this digit: done
                digits[i]++;
                return digits;
            }
            digits[i] = 0;                // 9 + 1 = 10: write 0, carry 1 to the left
        }
        digits.insert(digits.begin(), 1); // all digits were 9: 99...9 + 1 = 100...0
        return digits;
    }
};

int main() {                              // input: n, then the n digits
    int n;
    if (!(cin >> n)) return 0;
    vector<int> d(n);
    for (auto& x : d) cin >> x;
    vector<int> r = Solution().plusOne(d);
    for (size_t i = 0; i < r.size(); i++) cout << r[i] << (i + 1 < r.size() ? ' ' : '\n');
}
\end{lstlisting}
\begin{complexitybox}
\textbf{Time:} $\bigO{n}$ in the worst case ($99\ldots9$), $\bigO{1 + t}$ when the number ends with $t$ nines; averaged over random inputs, about $1.1$ digits are touched. \textbf{Space:} $\bigO{1}$ extra (the input is modified and returned; the \texttt{insert} is $\bigO{n}$ and happens at most once). \textbf{Verdict:} \ac{}.
\end{complexitybox}
\end{algorithmbox}

\begin{invariantbox}[Why the loop computes $N + 1$]
\textbf{Invariant.} When the loop is about to examine index $i$, digits $i+1, \dots, n-1$ have been set to 0, digits $0..i$ are unchanged, and the value still to be added is exactly $10^{\,n-1-i}$ (a carry of 1 into position $i$). Then (current array value) $+$ (pending carry) $= N + 1$.

\textbf{Proof.} Initially $i = n-1$, nothing is changed, pending $= 10^0 = 1$. If $d[i] < 9$, adding the carry to $d[i]$ produces a digit $\le 9$ --- a valid digit array whose value is exactly $N + 1$, so returning is correct. If $d[i] = 9$, replacing it by 0 removes $9\cdot 10^{\,n-1-i}$; together with the pending $10^{\,n-1-i}$ that is a change of $-10 \cdot 10^{\,n-1-i}$, compensated by a pending carry of $10^{\,n-i}$ into position $i-1$: the invariant holds for $i-1$. If the loop ends ($i = -1$), all digits are 0 and the pending carry is $10^n$: the answer is $1$ followed by $n$ zeros, which is what \texttt{insert} produces. \hfill$\blacksquare$
\end{invariantbox}

\begin{alternativebox}
\textbf{General addition with a carry variable.} Treat ``$+1$'' as the incoming carry and run the full school addition, building the result \emph{least significant digit first} with \texttt{push\_back}, then reverse. It never stops early, but the very same loop adds any digit array or any small integer $k$ (variant V1) --- it is the template to remember.
\begin{lstlisting}[style=cpp]
#include <bits/stdc++.h>
using namespace std;

class Solution {
public:
    vector<int> plusOne(vector<int>& digits) {
        vector<int> res;                       // digits of the result, least significant first
        int carry = 1;                         // "+1" enters as the initial carry
        for (int i = (int)digits.size() - 1; i >= 0; i--) {
            int s = digits[i] + carry;
            res.push_back(s % 10);
            carry = s / 10;
        }
        if (carry) res.push_back(carry);
        reverse(res.begin(), res.end());       // back to most significant first
        return res;
    }
};

int main() {                                   // input: n, then the n digits
    int n;
    if (!(cin >> n)) return 0;
    vector<int> d(n);
    for (auto& x : d) cin >> x;
    vector<int> r = Solution().plusOne(d);
    for (size_t i = 0; i < r.size(); i++) cout << r[i] << (i + 1 < r.size() ? ' ' : '\n');
}
\end{lstlisting}
Always $\bigO{n}$ time and $\bigO{n}$ extra memory for the new array. Growing at the \emph{end} (\texttt{push\_back}) instead of the front is why big-number libraries store digits least significant first.

\begin{diagrambox}
\begin{tikzpicture}[x=1cm, y=1cm]
    \dparray{g}{0.8}{input}{1/acell, 2/acell, 9/acell, 9/acell}
    \dparray{h}{-0.4}{\texttt{res}}{0/awrite, 0/awrite, 3/awrite, 1/awrite}
    \foreach \x/\y in {3/0, 2/1, 1/2, 0/3}{ \draw[faded] (g\x.south) -- (h\y.north); }
    \draw[flow] (4.0,0.2) -- node[lbl, above]{reverse} (5.4,0.2);
    \begin{scope}[xshift=5.7cm]
        \dparray{k}{0.2}{}{1/akeep, 3/akeep, 0/akeep, 0/akeep}
    \end{scope}
\end{tikzpicture}
\figcap{$[1,2,9,9]$: \texttt{res} collects $0, 0, 3, 1$ (least significant first, carry 0 at the end); reversing gives $[1,3,0,0]$.}
\end{diagrambox}
\end{alternativebox}

\begin{tricksbox}
\begin{itemize}[leftmargin=*]
    \item \textbf{Build the all-9s answer directly:} \texttt{vector<int> r(n + 1, 0); r[0] = 1;} instead of \texttt{insert} --- every digit is 0 anyway.
    \item \textbf{Count trailing 9s first:} $t$ = number of trailing 9s; if $t = n$ answer $10^n$, else set the last $t$ digits to 0 and increment $d[n-1-t]$. Same algorithm, two clearly separated phases.
    \item \textbf{No \texttt{long long} anywhere}: the only arithmetic is $d + 1 \le 10$.
    \item \textbf{Little-endian storage} (least significant digit first) makes growth free: the carry out of the top digit is a \texttt{push\_back}, and digit $i$ has weight $10^i$.
\end{itemize}
\begin{diagrambox}
\begin{tikzpicture}[x=1cm, y=1cm]
    \begin{scope}
        \node[font=\scriptsize\bfseries, text=deepblue, anchor=west] at (-0.6,1.0) {LeetCode order (big-endian)};
        \dparray{be}{0.3}{}{9/acell, 9/acell, 9/acell}
        \dparray{bf}{-0.9}{}{1/awrite, 0/awrite, 0/awrite, 0/awrite}
        \node[note] at (1.8,-1.75) {\texttt{insert(begin(), 1)}: every digit shifts, $\bigO{n}$};
    \end{scope}
    \begin{scope}[xshift=7.0cm]
        \node[font=\scriptsize\bfseries, text=deepblue, anchor=west] at (-0.6,1.0) {reversed order (little-endian)};
        \dparray{le}{0.3}{}{9/acell, 9/acell, 9/acell}
        \dparray{lf}{-0.9}{}{0/awrite, 0/awrite, 0/awrite, 1/akeep}
        \node[note] at (1.8,-1.75) {\texttt{push\_back(1)}: amortised $\bigO{1}$};
    \end{scope}
\end{tikzpicture}
\figcap{$999 + 1$ in both storage orders. The value is the same; only the position of the new digit differs.}
\end{diagrambox}
\end{tricksbox}

\subsubsection*{5. Trace}

\begin{tracebox}[{$[1,2,9,9]$: a carry through two 9s}]
\centering
\begin{tikzpicture}[x=1cm, y=1cm]
    \dparray{s0}{2.6}{start}{1/acell, 2/acell, 9/acell, 9/aread}
    \aptr{s03}{i=3}{accentorange}
    \dparray{s1}{0.9}{$i=3$}{1/acell, 2/acell, 9/aread, 0/awrite}
    \aptr{s12}{i=2}{accentorange}
    \dparray{s2}{-0.8}{$i=2$}{1/acell, 2/aread, 0/awrite, 0/awrite}
    \aptr{s21}{i=1}{accentorange}
    \dparray{s3}{-2.5}{$i=1$}{1/akeep, 3/akeep, 0/akeep, 0/akeep}
    \node[note, anchor=west, text width=5cm, align=left] at (4.2,0.9) {$d[3] = 9$: write 0, carry on};
    \node[note, anchor=west, text width=5cm, align=left] at (4.2,-0.8) {$d[2] = 9$: write 0, carry on};
    \node[note, anchor=west, text width=5cm, align=left] at (4.2,-2.5) {$d[1] = 2 < 9$: becomes 3, return $[1,3,0,0]$ --- $d[0]$ is never read};
\end{tikzpicture}
\figcap{Each row shows the array after the step named on the left; the orange pointer is the next index to examine.}

\medskip
\small
\begin{tabular}{@{}ccccl@{}}
\toprule
$i$ & \texttt{digits[i]} & case & carry after & \textbf{array} \\
\midrule
3 & 9 & $= 9$: write 0 & 1 & $[1,2,9,0]$ \\
2 & 9 & $= 9$: write 0 & 1 & $[1,2,0,0]$ \\
1 & 2 & $< 9$: increment, return & 0 & $[1,3,0,0]$ \\
\bottomrule
\end{tabular}
\end{tracebox}

\begin{tracebox}[{$[9,9,9]$: the array grows}]
\centering
\begin{tikzpicture}[x=1cm, y=1cm]
    \dparray{u0}{1.2}{start}{{}/emptycell, 9/acell, 9/acell, 9/acell}
    \dparray{u1}{0}{after loop}{{}/emptycell, 0/awrite, 0/awrite, 0/awrite}
    \dparray{u2}{-1.2}{insert 1}{1/akeep, 0/akeep, 0/akeep, 0/akeep}
    \node[note, anchor=west, text width=5.2cm, align=left] at (4.3,0) {$i = 2, 1, 0$: three 9s become 0; the loop ends with a pending carry of $10^3$, which becomes the new leading digit: $999 + 1 = 1000$};
\end{tikzpicture}
\figcap{The only input shape that changes the length. The empty grey cell is the position created by \texttt{insert}.}

\medskip
\small
\begin{tabular}{@{}ccccl@{}}
\toprule
$i$ & \texttt{digits[i]} & case & carry after & \textbf{array} \\
\midrule
2 & 9 & write 0 & 1 & $[9,9,0]$ \\
1 & 9 & write 0 & 1 & $[9,0,0]$ \\
0 & 9 & write 0 & 1 & $[0,0,0]$ \\
--- & --- & loop ended: insert 1 & 0 & $[1,0,0,0]$ \\
\bottomrule
\end{tabular}
\end{tracebox}

\subsubsection*{6. Test suite}

\begin{testbox}
\centering\small
\begin{tabular}{@{}lll@{}}
\toprule
\textbf{digits} & \textbf{Why it is interesting} & \textbf{Expected} \\
\midrule
$[1,2,3]$, $[4,3,2,1]$ & official Samples 1--2: no carry & $[1,2,4]$, $[4,3,2,2]$ \\
$[9]$ & official Sample 3: growth & $[1,0]$ \\
$[0]$ & the number zero (the only leading zero allowed) & $[1]$ \\
$[1,2,9,9]$ & carry through a run of 9s & $[1,3,0,0]$ \\
$[8,9,9,9]$ & carry stops at the top digit & $[9,0,0,0]$ \\
$[9,9,9]$ & all 9s: length $n+1$ & $[1,0,0,0]$ \\
twenty 9s & beyond 64 bits: the conversion approach fails & $1$ and twenty 0s \\
$[9,0,9]$ & a 9 that is \emph{not} part of the trailing run & $[9,1,0]$ \\
\bottomrule
\end{tabular}

\medskip
The digit loop and the carry-variable version reproduce the official samples and agree with Python's big integers on 500 random digit arrays of length up to 25 (biased towards 9s).
\end{testbox}

\subsubsection*{7. Common mistakes}

\begin{warningbox}[Mistake 1 --- converting to an integer]
\texttt{long long} holds 18 full digits, \texttt{unsigned long long} 19; the input has up to 100. See the brute force.
\end{warningbox}

\begin{warningbox}[Mistake 2 --- forgetting the all-9s case]
\begin{lstlisting}[style=snippet]
for (i = n - 1; i >= 0; i--) { ... }  return digits;       // WRONG: [9] -> [0]
... digits.insert(digits.begin(), 1); return digits;       // RIGHT
\end{lstlisting}
\end{warningbox}

\begin{warningbox}[Mistake 3 --- not stopping after the increment]
Without the \texttt{return} inside the loop, the next iteration sees a digit $< 9$ again and increments it too: $[1,2,3] \to [2,3,4]$. The carry must be consumed exactly once.
\end{warningbox}

\begin{warningbox}[Mistake 4 --- \texttt{digits.size() - 1} with an unsigned index]
\texttt{for (size\_t i = digits.size() - 1; i >= 0; i--)} never ends: an unsigned \texttt{i} is always $\ge 0$ and wraps from 0 to $2^{64}-1$. Use \texttt{int i = (int)digits.size() - 1}.
\end{warningbox}

\subsubsection*{8. Takeaways}

\begin{itemize}[leftmargin=*]
    \item Numbers longer than 18 digits are digit arrays; operate on the digits, never convert.
    \item The carry of an addition with two digits is 0 or 1; for $+1$ it only travels through 9s, so the loop stops at the first non-9.
    \item The result can be one digit longer than the input --- the only border case, handled once at the end.
    \item Least-significant-first storage turns ``grow at the front'' into \texttt{push\_back}.
\end{itemize}

\begin{variantsbox}
\variant{Add to Array-Form of Integer}{\exeasy\ \lc{Easy}\ \varlink{Add to Array-Form of Integer}{LeetCode 989}{https://leetcode.com/problems/add-to-array-form-of-integer/}}
  {add an integer $k \le 10^4$ instead of 1; the result may be several digits longer.}
  {model: the carry is no longer 0/1: keep $k$ itself as the carry and add $d[i]$ to it.}
  {algorithm: \texttt{while (i >= 0 || k > 0) \{ if (i >= 0) k += d[i--]; push k \% 10; k /= 10; \}}, reverse. Samples $\to$ 1234, 455, 1021.}
\variant{Add Binary}{\exeasy\ \lc{Easy}\ \varlink{Add Binary}{LeetCode 67}{https://leetcode.com/problems/add-binary/}}
  {two numbers given as binary strings of length up to $10^4$.}
  {model: digit arrays in base 2; $s = a_i + b_i + c \le 3$.}
  {algorithm: the carry loop with \texttt{s \% 2} and \texttt{s / 2}; samples $\to$ \texttt{100}, \texttt{10101}.}
\variant{Add Two Numbers}{\exmedium\ \lc{Medium}\ \varlink{Add Two Numbers}{LeetCode 2}{https://leetcode.com/problems/add-two-numbers/}}
  {the two numbers are linked lists of digits, least significant first.}
  {model: the alternative's carry loop, with ``next node'' instead of ``index minus one''.}
  {algorithm: walk both lists while either remains or the carry is 1; samples $\to [7,0,8]$, $[0]$, $[8,9,9,9,0,0,0,1]$.}
\variant{Minus One}{\varnone}
  {return the digits of $N - 1$ for $N \ge 1$.}
  {model: the borrow travels through 0s, which become 9s; the length may \emph{shrink} ($1000 \to 999$).}
  {algorithm: mirror loop (stop at the first non-zero digit, decrement it), then remove a leading zero if $n > 1$.}
\end{variantsbox}
